% Report-ready section for the Network Flow component.
\section{Water Distribution as a Maximum-Flow Problem}
\label{sec:water-flow}

\subsection{Problem abstraction}
Let the physical water-distribution system be a directed graph
$G=(V,E)$, where each pipe $(u,v)\in E$ has a nonnegative capacity
$c(u,v)$. Let $P\subseteq V$ be supply nodes, with available supply
$s(p)\ge 0$, and let $D\subseteq V$ be demand nodes, with demand cap
$d(z)\ge 0$. The goal is to maximize the total water delivered to demand
nodes while respecting pipe capacities, supply limits, and demand limits.
This abstraction models capacity-constrained allocation; it does not model
hydraulic pressure, head loss, pipe dynamics, or water quality.

\subsection{Reduction construction}
Construct a flow network $G'=(V',E')$ by introducing a super-source $S$
and a super-sink $T$. Set $V'=V\cup\{S,T\}$. Retain every physical
pipe $(u,v)\in E$ with capacity $c(u,v)$. For each supply node $p\in P$,
add an edge $(S,p)$ with capacity $s(p)$. For each demand node $z\in D$,
add an edge $(z,T)$ with capacity $d(z)$. Run a maximum-flow algorithm
from $S$ to $T$ in $G'$.

\subsection{Correctness proof}
We show that feasible allocations in the original abstraction correspond
to feasible flows in $G'$ and that the objective value is preserved.

\paragraph{From a water allocation to a flow.}
Consider any feasible allocation of water through the physical network.
Assign the amount dispatched from each supply node $p$ to edge $(S,p)$.
For every physical pipe $(u,v)$, assign the amount of water sent through
that pipe to its corresponding edge in $G'$. Finally, assign the amount
received at demand node $z$ to edge $(z,T)$. Supply limits, pipe capacities,
and demand caps ensure that all edge-capacity constraints are satisfied.
At every intermediate node, conservation of water is exactly the flow
conservation constraint. Thus this assignment is a feasible $S$--$T$ flow.
Its value equals the total water delivered to the demand nodes.

\paragraph{From a flow to a water allocation.}
Conversely, take any feasible $S$--$T$ flow in $G'$. Interpret the flow
on each physical-pipe edge as the amount transported through that pipe,
the flow on $(S,p)$ as the amount drawn from supply node $p$, and the
flow on $(z,T)$ as the amount delivered to demand node $z$. Capacity
constraints enforce each pipe capacity, supply limit, and demand cap.
Flow conservation at intermediate vertices ensures that transported water
is balanced. The resulting allocation is therefore feasible in the original
abstraction, and its total delivered amount equals the flow value.

\paragraph{Optimality.}
The two mappings above preserve feasibility and objective value. Hence the
maximum possible water delivery in the abstraction is equal to the maximum
$S$--$T$ flow value in $G'$. A maximum-flow algorithm therefore produces
an optimal allocation for this model.

\subsection{Construction and algorithm complexity}
The transformation adds two vertices and at most $|P|+|D|$ edges beyond
the physical network. Its construction takes
$O(|V|+|E|+|P|+|D|)$ time under an adjacency-list representation.
The implementation then computes a maximum flow using Dinic's algorithm.
For a general directed network, Dinic's algorithm has the standard
$O(|V'|^2|E'|)$ worst-case bound. Since $|V'|=|V|+2$ and
$|E'|=|E|+|P|+|D|$, this is
$O((|V|+2)^2(|E|+|P|+|D|))$ for the constructed network.